\chapter{Forces, Phase Space, and Conservation Laws}

\lectureoverview{A law of motion must do more than relate familiar physical
quantities.  It must specify the state of a system, carry that state from one
instant to the next, and reveal which distinctions and quantities survive the
motion.  We shall reach Newtonian mechanics by first testing a plausible but
wrong force law, then enlarge the state from position to phase space, and
finally derive conservation of momentum and energy.}

\section{A Useful Wrong Law}

Everyday motion is dominated by friction.  Push a heavy object across a floor
and it moves; stop pushing and it soon stops.  A larger push generally produces
a larger speed, while a heavier object requires a harder push.  In such a
world, it is natural to suppose that force produces velocity itself.  This was
the basic intuition behind the Aristotelian description of motion.

Ice skating supplies the missing experience.  A skater needs force to start,
stop, or change direction, but not to continue gliding once in motion.  Had
Aristotle worked on nearly frictionless ice rather than ordinary ground, the
defect in the force--velocity intuition would have been much harder to miss.

To expose what is wrong with that intuition, let us give it a precise modern
form.  We write the deliberately incorrect equation

\begin{equation}
  \mathbf F=m\mathbf v.
  \label{eq:aristotle-vector-law}
\end{equation}

This is not Newton's law and should not be memorized as physics.  Its value is
diagnostic: once a verbal rule has become an equation, we can ask whether it
predicts motion and whether it preserves enough information to reconstruct the
past.  The brief classroom discussion of the Greek word for force points to
\emph{dynamis}, the root also heard in words such as dynamo and dynamics, but
the terminology is secondary to the test of the law itself.

\begin{figure}[htbp]
  \centering
  \includegraphics[width=0.82\textwidth]{lecture_02_figure_02.png}
  \caption{The intentionally wrong force--velocity law used to test
  predictability and reversibility.}
  \lectureframe{00:04:51}
\end{figure}

Restrict the motion to one spatial coordinate.  Since \(v=\dot x\), the law is

\begin{equation}
  F(t)=m\dot x(t).
  \label{eq:aristotle-1d}
\end{equation}

Now divide time into short intervals of length \(\delta\).  Over one interval,

\begin{equation}
  \dot x(t)\simeq
  \frac{x(t+\delta)-x(t)}{\delta}.
  \label{eq:first-difference}
\end{equation}

Substitution into Eq.~\eqref{eq:aristotle-1d} gives

\begin{equation}
  x(t+\delta)
  =x(t)+\frac{\delta}{m}F(t).
  \label{eq:aristotle-update}
\end{equation}

This has the form of a genuine update rule.  If the present position and the
force are known, the position at the next stroboscopic instant follows.  When
the force vanishes, however, Aristotle's rule gives \(\dot x=0\): the object
must remain at rest.  Motion without a sustaining force is excluded from the
start.

\begin{figure}[htbp]
  \centering
  \includegraphics[width=0.84\textwidth]{lecture_02_figure_07.png}
  \caption{The discrete Aristotle update with a restoring force.  Each step
  multiplies the displacement by a factor slightly smaller than one.}
  \lectureframe{00:10:15}
\end{figure}

\subsection{A spring in Aristotle's mechanics}

Consider a particle attached to a spring, with equilibrium at \(x=0\).  A
displacement produces the restoring force

\begin{equation}
  F(x)=-kx,
  \qquad k>0.
  \label{eq:hooke-law}
\end{equation}

Inserting this force into Eq.~\eqref{eq:aristotle-update} yields

\begin{equation}
  x(t+\delta)
  =\left(1-\frac{k\delta}{m}\right)x(t).
  \label{eq:aristotle-spring-update}
\end{equation}

For sufficiently small \(\delta\), the factor in parentheses lies just below
one.  Every update therefore reduces the displacement by the same fraction:

\begin{equation}
  x_n=\left(1-\frac{k\delta}{m}\right)^n x_0.
  \label{eq:aristotle-geometric-decay}
\end{equation}

The continuous equation and its solution are

\begin{align}
  m\dot x&=-kx,
  \label{eq:aristotle-spring-continuous}\\
  x(t)&=x(0)e^{-kt/m}.
  \label{eq:aristotle-spring-solution}
\end{align}

The spring does not oscillate.  Released from any displacement, it approaches
the origin exponentially, moving more and more slowly as it comes in.

\begin{figure}[htbp]
  \centering
  \includegraphics[width=0.84\textwidth]{lecture_02_figure_08.png}
  \caption{The exponential solution of the Aristotelian spring equation.  The
  displacement decays rather than oscillates.}
  \lectureframe{00:12:02}
\end{figure}

The future is easy to predict: every initial displacement approaches the same
resting configuration.  The past is another matter.  Suppose the particle is
observed very near the origin with only finite accuracy.  It could have begun
at any of many distinct displacements, because all of their exponentially
small remnants eventually lie inside the same measurement tolerance.  With
perfect mathematical precision Eq.~\eqref{eq:aristotle-spring-solution} can be
run backward, but practical distinctions between histories are steadily
erased.  The law is therefore suspicious from the standpoint of reversible
fundamental dynamics.

This does not mean that Aristotle was irrational.  The friction-dominated
motion he described really is dissipative.  The mistake was to promote that
effective behavior to a fundamental law.  Friction is itself a force arising
from complicated microscopic interactions.  When a hand pushes an object at
roughly steady speed, the applied force is balancing friction; the applied
force is not creating velocity as such.

\section{Newton's Law and Its Physical Meaning}

The Newtonian correction is that force changes motion:

\begin{equation}
  \mathbf F=m\mathbf a=m\ddot{\mathbf x}.
  \label{eq:newton-second-law}
\end{equation}

Before treating this as an evolution equation, two apparent circularities must
be untangled.  The first concerns frames of reference.  One often hears that
Newton's laws hold in inertial frames and that an inertial frame is one in
which Newton's laws hold.  The physical assertion is not that circular
definition.  It is the empirical claim that there \emph{exist} frames,
described by suitable rods and clocks, in which Newton's laws take their simple
form.  Frames accelerating or rotating relative to them require additional
apparent forces.  How the equations transform in those other frames can be
worked out later.

The second circularity concerns force and mass.  If force is defined as mass
times acceleration and mass as force divided by acceleration, the equation
seems empty.  An operational experiment separates the definitions from the
law.  Compare an object with a standard kilogram by balancing the two, then
attach the object to a spring balance in otherwise force-free space.  Define
one unit of force by a fixed small extension of the spring.  Two identical
springs, each stretched by that same amount, provide two units of force and
produce twice the acceleration.  It is safer to add a second spring than to
stretch one spring twice as far, because a real spring eventually becomes
nonlinear or breaks.  Keeping one spring but joining two equal masses produces
half the acceleration; three equal masses produce one third.

\begin{figure}[htbp]
  \centering
  \includegraphics[width=0.84\textwidth]{lecture_02_figure_09.png}
  \caption{The spring-balance thought experiment used to define equal units of
  force and test acceleration proportional to force.}
  \lectureframe{00:22:45}
\end{figure}

These comparisons give empirical content to

\begin{equation}
  a\propto F,
  \qquad
  a\propto\frac{1}{m},
  \qquad
  F=ma.
  \label{eq:force-mass-scaling}
\end{equation}

The word \emph{particle} will be used for any small object whose internal
structure is irrelevant, not only for an elementary particle.

\subsection{Why position alone is not a state}

We now test Newton's equation by the same stroboscopic method.  The centered
finite difference for acceleration is

\begin{equation}
  \ddot x(t)\simeq
  \frac{x(t+\delta)-2x(t)+x(t-\delta)}{\delta^2}.
  \label{eq:second-difference}
\end{equation}

Thus Newton's equation implies

\begin{equation}
  x(t+\delta)
  =2x(t)-x(t-\delta)+\frac{\delta^2}{m}F(t).
  \label{eq:newton-position-update}
\end{equation}

The next position cannot be obtained from \(x(t)\) alone.  We also need the
previous position \(x(t-\delta)\), or equivalently the velocity at the present
instant.  At first this makes Newton's equation appear less predictive than
Aristotle's.  The correct conclusion is instead that position alone was an
incomplete description of the state.  To know everything relevant at one
instant means to know enough data to determine the subsequent motion.  For a
Newtonian particle those data are position and velocity.

\section{Momentum and Phase Space}

Introduce momentum by

\begin{equation}
  \mathbf p=m\mathbf v=m\dot{\mathbf x}.
  \label{eq:momentum-definition}
\end{equation}

The mass of a fixed Newtonian system is taken to be constant.  An object can
lose a piece, as a pen can lose its cap, but then the object being described
has changed; the chosen closed system still has constant total mass.  With
constant \(m\), Newton's second law becomes

\begin{equation}
  \mathbf F=\dot{\mathbf p}.
  \label{eq:newton-momentum-form}
\end{equation}

\begin{figure}[htbp]
  \centering
  \includegraphics[width=0.82\textwidth]{lecture_02_figure_03.png}
  \caption{Newton's law written as the time derivative of momentum.}
  \lectureframe{00:32:35}
\end{figure}

For one coordinate, Newtonian motion can now be written as two first-order
equations for two independent state variables:

\begin{align}
  \dot p&=F(x,t),
  \label{eq:first-order-p}\\
  \dot x&=\frac{p}{m}.
  \label{eq:first-order-x}
\end{align}

One equation says how the force changes momentum; the other says how momentum
changes position.  Their first-order discrete form is

\begin{align}
  p(t+\delta)&=p(t)+\delta F(x(t),t)+O(\delta^2),
  \label{eq:momentum-discrete}\\
  x(t+\delta)&=x(t)+\frac{\delta}{m}p(t)+O(\delta^2).
  \label{eq:position-discrete}
\end{align}

Given \(x(t)\) and \(p(t)\), these equations determine the next pair, which
then becomes the input for the next step.  The trajectory is built by repeated
updates.

\begin{figure}[htbp]
  \centering
  \includegraphics[width=0.84\textwidth]{lecture_02_figure_04.png}
  \caption{The two first-order equations and the corresponding initial data.
  Position and momentum together specify the state.}
  \lectureframe{00:35:47}
\end{figure}

The space of all pairs \((x,p)\) is called \emph{phase space}.  For one
particle moving in three dimensions it has six coordinates,

\begin{equation}
  (x,y,z,p_x,p_y,p_z).
  \label{eq:one-particle-phase-space}
\end{equation}

Phase space is the space of complete initial conditions.  Configuration space
records positions only; phase space records precisely the data needed by the
law of motion.

\begin{classroomqa}
  \audiencequestion{If the force is a function of position rather than a known
  function of time, does the update become more complicated?\lecturetimestamp{00:36:27--00:37:07}}
  \lecturerresponse{The equations remain predictive.  Once \(x\) and \(p\) are
  known, the value \(F(x)\) is known.  Equation~\eqref{eq:first-order-p}
  advances \(p\), Eq.~\eqref{eq:first-order-x} advances \(x\), and the updated
  pair can be used again.}
\end{classroomqa}

\section{The Newtonian Spring}

Return to the spring so that the two laws can be compared on identical
physical ground.  Newton gives

\begin{equation}
  m\ddot x=-kx.
  \label{eq:newton-spring}
\end{equation}

Set \(m=k=1\) temporarily to remove constants.  Then

\begin{equation}
  \ddot x=-x.
  \label{eq:normalized-oscillator}
\end{equation}

Unlike the exponential function, a sine or cosine returns to minus itself
after two differentiations.  The general solution may be written

\begin{align}
  x(t)&=C\cos(t-t_0),
  \label{eq:oscillator-x}\\
  p(t)=\dot x(t)&=-C\sin(t-t_0).
  \label{eq:oscillator-p}
\end{align}

The two constants carry the two pieces of initial information.  The amplitude
\(C\) gives the maximum displacement, while \(t_0\) fixes the phase, chosen
here so that the velocity vanishes at \(t=t_0\).  The particle stops only
momentarily at the endpoints.  It then returns through the origin, where its
speed and momentum are largest, and continues oscillating.

Squaring Eqs.~\eqref{eq:oscillator-x} and \eqref{eq:oscillator-p} gives

\begin{equation}
  x^2+p^2=C^2.
  \label{eq:oscillator-phase-circle}
\end{equation}

Thus the phase-space point moves on a circle of radius \(C\).  Different
initial amplitudes produce different circles; changing \(t_0\) chooses a
different starting point on the same circle.  Every circle is traversed with
the same period \(2\pi\) in these normalized units, independent of amplitude.

\begin{figure}[htbp]
  \centering
  \includegraphics[width=0.84\textwidth]{lecture_02_figure_10.png}
  \caption{The oscillator solution and its circular orbit in the \((x,p)\)
  phase plane.}
  \lectureframe{00:43:20}
\end{figure}

\begin{figure}[htbp]
  \centering
  \begin{tikzpicture}[scale=1.05,>=Latex]
    \draw[->] (-2.6,0) -- (2.75,0) node[right] {$x$};
    \draw[->] (0,-2.4) -- (0,2.55) node[above] {$p$};
    \draw[blue!65!black,thick] (0,0) circle (1.85);
    \draw[->,blue!65!black,thick]
      (1.60,0.93) arc[start angle=30,end angle=72,radius=1.85];
    \fill[red!65!black] (1.60,0.93) circle (2pt);
    \draw[dashed] (0,0) -- (1.60,0.93)
      node[midway,below right] {$C$};
    \node[align=left] at (-1.55,1.85)
      {$x=C\cos(t-t_0)$\\$p=-C\sin(t-t_0)$};
  \end{tikzpicture}
  \caption{A phase-space orbit retains both position and momentum.  Its radius
  labels the motion, while the angle advances uniformly in time.}
  \label{fig:oscillator-phase-space}
\end{figure}

Projection onto the horizontal axis gives the familiar back-and-forth motion
of \(x\); projection onto the vertical axis gives the oscillation of \(p\).
Unlike the Aristotelian spring, distinct circles never collapse into one
another.  Knowing a phase-space point determines both its future and its past.

For general \(m\) and \(k\), the angular frequency is

\begin{equation}
  \omega=\sqrt{\frac{k}{m}},
  \qquad
  x(t)=C\cos\bigl(\omega(t-t_0)\bigr),
  \label{eq:general-oscillator-frequency}
\end{equation}

and the phase-space orbit becomes an ellipse in the unscaled \((x,p)\) plane.
The circular picture is the convenient \(m=k=1\) normalization used in the
lecture.

\begin{classroomqa}
  \audiencequestion{If the second derivative in the oscillator equation were
  replaced by a first derivative, would predictability be lost?\lecturetimestamp{00:45:23--00:46:04}}
  \lecturerresponse{No.  A first-order equation can still predict forward
  motion from suitable initial data.  The stronger issue is whether its flow
  preserves distinctions well enough to be reversible; the relevant general
  theorem is deferred until the needed machinery is available.}
\end{classroomqa}

\section{Time Reversal}

The oscillator picture makes reversibility concrete, but there is also a
direct test on the equations.  Introduce a reversed time coordinate

\begin{equation}
  T=-t.
  \label{eq:reverse-time}
\end{equation}

Then

\begin{equation}
  \frac{d}{dt}=-\frac{d}{dT},
  \qquad
  \frac{d^2}{dt^2}=\frac{d^2}{dT^2}.
  \label{eq:time-derivative-parity}
\end{equation}

A first derivative changes sign under time reversal, while a second derivative
does not.  Consequently the autonomous Aristotelian equation

\begin{equation}
  m\dot x=F(x)
  \label{eq:aristotle-time-reversal}
\end{equation}

does not retain its form when a movie is run backward.  A decaying spring would
become a spring that spontaneously departs from equilibrium.  By contrast,

\begin{equation}
  m\ddot x=F(x)
  \label{eq:newton-time-reversal}
\end{equation}

is unchanged.  Reversing a Newtonian movie reverses velocity but leaves
acceleration unchanged, so the backward movie is another allowed solution.

\begin{figure}[htbp]
  \centering
  \includegraphics[width=0.84\textwidth]{lecture_02_figure_11.png}
  \caption{The board comparison under \(t\mapsto -t\): first derivatives
  reverse sign, whereas second derivatives do not.}
  \lectureframe{00:50:45}
\end{figure}

\begin{classroomqa}
  \audiencequestion{Why should reversibility be accepted as a premise about
  physics in general?\lecturetimestamp{00:46:16--00:48:26}}
  \lecturerresponse{It is tested rather than guaranteed by definition.  We
  examine successful laws, identify common experimental features, and elevate
  those features to principles.  The deeper logical order runs in the opposite
  direction: quantum field theory is more fundamental, and Newtonian mechanics
  emerges through approximations such as slow motion and heavy objects.  A
  course must build upward from simpler laws, so reversibility is introduced
  now as a principle whose reach will be tested later.}
\end{classroomqa}

\begin{classroomqa}
  \audiencequestion{Was reversibility itself a problem Newton consciously
  recognized?\lecturetimestamp{00:48:28--00:49:02}}
  \lecturerresponse{The historical answer is uncertain.  Newton certainly
  understood inertia: force is required to change a state of motion, not to
  sustain uniform motion.  Whether he framed that insight explicitly in terms
  of reversibility is a separate historical question.}
\end{classroomqa}

\section{Cycles and Conserved Quantities}

The phase-space circles suggest the next idea.  A trajectory remains on one
circle and never jumps to another.  Therefore the circle can be labeled by a
number that remains constant along the motion.  For the normalized oscillator
that number is

\begin{equation}
  C^2=x^2+p^2.
  \label{eq:cycle-label}
\end{equation}

Up to a factor of one half, this is the oscillator's energy.  Larger circles
have larger energy; changing the starting phase does not change it.  We shall
derive that statement rather than merely read it from the picture, but first
consider momentum.

Newton's first law says that an object subject to no net force moves with
constant velocity.  It is already contained in the second law:

\begin{equation}
  \mathbf F=0
  \quad\Longrightarrow\quad
  \dot{\mathbf p}=0
  \quad\Longrightarrow\quad
  \mathbf v=\text{constant}
  \label{eq:newton-first-from-second}
\end{equation}

for constant mass.  Newton's third law contributes genuinely new information
about interactions between particles.

\subsection{Equal, opposite, and central forces}

Let \(\mathbf F_{ij}\) denote the force on particle \(i\) exerted by particle
\(j\).  The equal-and-opposite part of the third law is

\begin{equation}
  \mathbf F_{ij}=-\mathbf F_{ji}.
  \label{eq:newton-third-pair}
\end{equation}

If the two forces pointed in the same direction, an isolated pair could
accelerate spontaneously.  It would be a propellant-free machine producing
motion from internal forces alone.  Such behavior is not observed.

\begin{figure}[htbp]
  \centering
  \includegraphics[width=0.84\textwidth]{lecture_02_figure_12.png}
  \caption{The indexed statement of equal and opposite interaction forces.}
  \lectureframe{00:59:25}
\end{figure}

For the point-particle forces considered here, the pair also acts along the
line joining the particle centers.  Equal and opposite forces that were not
collinear could form a pure torque and set an isolated rigid pair rotating.
The central-force condition prevents that possibility:

\begin{equation}
  \mathbf F_{ij}\parallel
  (\mathbf r_i-\mathbf r_j).
  \label{eq:central-force-condition}
\end{equation}

The equal-and-opposite condition will give conservation of linear momentum.
The line-of-centers condition is tied to conservation of angular momentum,
which will be studied later.

\begin{figure}[htbp]
  \centering
  \includegraphics[width=0.84\textwidth]{lecture_02_figure_13.png}
  \caption{The central-force part of Newton's third law: interaction forces
  lie along the line joining the particles.}
  \lectureframe{01:01:38}
\end{figure}

\subsection{Total momentum}

Consider a closed system of \(N\) particles.  There are no external forces;
the only forces are those exerted by the particles on one another.  Force
additivity and Newton's second law give, for each particle,

\begin{equation}
  \dot{\mathbf p}_i
  =m_i\ddot{\mathbf r}_i
  =\sum_{\substack{j=1\\j\ne i}}^N \mathbf F_{ij}.
  \label{eq:many-particle-newton}
\end{equation}

Sum this equation over \(i\):

\begin{align}
  \frac{d}{dt}\sum_{i=1}^N\mathbf p_i
  &=\sum_{i=1}^N
    \sum_{\substack{j=1\\j\ne i}}^N\mathbf F_{ij}
  \notag\\
  &=\sum_{i<j}\bigl(\mathbf F_{ij}+\mathbf F_{ji}\bigr)
  =0.
  \label{eq:momentum-pair-cancellation}
\end{align}

Every interaction appears twice, once in each direction, and the pair cancels.
Defining

\begin{equation}
  \mathbf P_{\mathrm{tot}}=\sum_{i=1}^N\mathbf p_i,
  \label{eq:total-momentum}
\end{equation}

we obtain the first conservation law:

\begin{equation}
  \frac{d\mathbf P_{\mathrm{tot}}}{dt}=0.
  \label{eq:momentum-conservation}
\end{equation}

\begin{figure}[htbp]
  \centering
  \includegraphics[width=0.84\textwidth]{lecture_02_figure_14.png}
  \caption{Summing the particle equations makes the pairwise forces cancel,
  leaving the derivative of total momentum equal to zero.}
  \lectureframe{01:04:38}
\end{figure}

Historically we have presented momentum conservation as a consequence of
Newton's third law.  Modern physics often regards momentum conservation as the
more fundamental statement: if an isolated two-particle system conserves total
momentum, its internal momentum changes must be equal and opposite.  For the
present development, however, the Newtonian route is the useful one.

\section{Energy in One Dimension}

Take one particle moving on a line under a force that depends only on its
position:

\begin{equation}
  F=F(x).
  \label{eq:position-force}
\end{equation}

Such a force is called time independent.  This does not mean that the force
experienced by the moving particle has a constant numerical value.  Its value
changes as \(x(t)\) changes.  It means that whenever the particle is at the
same position \(x\), it experiences the same force, regardless of the clock
time.

In one dimension, any sufficiently regular force function can be integrated.
We therefore define a potential energy \(V(x)\) by

\begin{equation}
  F(x)=-\frac{dV}{dx},
  \label{eq:potential-definition-1d}
\end{equation}

or equivalently

\begin{equation}
  V(x)=-\int^x F(u)\,du+V_0.
  \label{eq:potential-integral}
\end{equation}

The additive constant \(V_0\) has no effect on the force.  Picture \(V(x)\) as
the height of a landscape above the \(x\)-axis.  The particle is pushed toward
lower height, and the force is larger where the slope is steeper.  At the top
or bottom of a smooth hill the slope vanishes, so the force vanishes there.

\begin{figure}[htbp]
  \centering
  \includegraphics[width=0.84\textwidth]{lecture_02_figure_15.png}
  \caption{Potential energy as a one-dimensional landscape.  The force points
  downhill and has magnitude set by the slope.}
  \lectureframe{01:12:57}
\end{figure}

Define kinetic, potential, and total energy by

\begin{equation}
  T=\frac12m\dot x^2,
  \qquad
  E=T+V(x)
  =\frac12m\dot x^2+V(x).
  \label{eq:energy-1d}
\end{equation}

At this stage energy is simply the name assigned to this candidate conserved
quantity.  Its physical interpretation is suggested by the landscape: rolling
downhill lowers \(V\) while increasing speed and therefore \(T\); climbing
uphill does the reverse.  Newton's equation proves that the exchange is exact.

Differentiate Eq.~\eqref{eq:energy-1d}.  Since \(m\) is constant and \(V\)
depends on time only through \(x(t)\),

\begin{align}
  \frac{dE}{dt}
  &=m\dot x\ddot x+\frac{dV}{dx}\dot x
  \notag\\
  &=\dot x\left(m\ddot x+\frac{dV}{dx}\right).
  \label{eq:energy-derivative-1d}
\end{align}

Newton's law and the definition of \(V\) say

\begin{equation}
  m\ddot x=F=-\frac{dV}{dx}.
  \label{eq:newton-potential-1d}
\end{equation}

The bracket in Eq.~\eqref{eq:energy-derivative-1d} therefore vanishes:

\begin{equation}
  \frac{dE}{dt}=0.
  \label{eq:energy-conservation-1d}
\end{equation}

\begin{figure}[htbp]
  \centering
  \includegraphics[width=0.84\textwidth]{lecture_02_figure_16.png}
  \caption{The one-dimensional energy derivative factored into velocity times
  Newton's equation.}
  \lectureframe{01:19:18}
\end{figure}

No additional dynamical assumption was needed in one dimension beyond a
regular position-dependent force, because every such force admits an
antiderivative.

\section{Several Dimensions and Conservative Forces}

For one particle with coordinates \(x_i\), Newton's equation has components

\begin{equation}
  F_i(\mathbf x)=m\ddot x_i.
  \label{eq:newton-components}
\end{equation}

The force component \(F_i\) may depend on all coordinates, not only on
\(x_i\).  We would like to find one scalar function \(V(\mathbf x)\) such that

\begin{equation}
  F_i(\mathbf x)
  =-\frac{\partial V(\mathbf x)}{\partial x_i},
  \qquad
  \mathbf F=-\boldsymbol\nabla V.
  \label{eq:gradient-force}
\end{equation}

Unlike the one-dimensional construction, this is not automatic.  Three
arbitrary functions \(F_1,F_2,F_3\) generally contain more independent
information than the derivatives of one scalar function.  Equivalently, a
potential exists only when the force field satisfies compatibility conditions
such as equality of the appropriate mixed partial derivatives; in a simply
connected region this is expressed by \(\boldsymbol\nabla\times\mathbf F=0\).
We therefore take Eq.~\eqref{eq:gradient-force} as a physical postulate for
the class of forces under consideration, not as a mathematical identity.

\begin{figure}[htbp]
  \centering
  \includegraphics[width=0.84\textwidth]{lecture_02_figure_17.png}
  \caption{The multidimensional force law and a contour sketch of the
  potential.  Force points across level curves toward decreasing \(V\).}
  \lectureframe{01:24:15}
\end{figure}

On a contour map, \(V\) is constant along each level curve.  The gradient is
perpendicular to those curves and points toward steepest increase; the force
\(-\boldsymbol\nabla V\) points toward steepest descent.

\begin{figure}[htbp]
  \centering
  \begin{tikzpicture}[scale=1.0,>=Latex]
    \draw[->] (-3.0,0) -- (3.1,0) node[right] {$x_1$};
    \draw[->] (0,-2.4) -- (0,2.55) node[above] {$x_2$};
    \draw[orange!65!black,thick]
      plot[smooth cycle] coordinates
      {(-0.55,-0.35) (0.15,-0.55) (0.72,-0.10) (0.50,0.53) (-0.25,0.65)};
    \draw[orange!65!black,thick]
      plot[smooth cycle] coordinates
      {(-1.15,-0.75) (0.20,-1.10) (1.35,-0.25) (1.10,0.95) (-0.45,1.25)};
    \draw[orange!65!black,thick]
      plot[smooth cycle] coordinates
      {(-1.85,-1.20) (0.15,-1.70) (2.15,-0.40) (1.70,1.50) (-0.80,1.85)};
    \draw[->,blue!65!black,very thick] (1.62,0.73) -- (0.82,0.20)
      node[midway,above left] {$\mathbf F$};
    \draw[dashed,->,red!65!black,thick] (0.82,0.20) -- (1.62,0.73)
      node[midway,below right] {$\boldsymbol\nabla V$};
    \node[orange!65!black] at (-1.62,1.45) {$V=V_3$};
    \node[orange!65!black] at (-0.92,0.92) {$V=V_2$};
    \node[orange!65!black] at (-0.35,0.30) {$V=V_1$};
  \end{tikzpicture}
  \caption{Level curves of a scalar potential.  The gradient is normal to the
  contours, and a conservative force points in the opposite direction.}
  \label{fig:potential-contours}
\end{figure}

When Eq.~\eqref{eq:gradient-force} holds, define

\begin{equation}
  E=\frac12m\sum_i\dot x_i^2+V(x_1,\ldots,x_n).
  \label{eq:energy-many-coordinates}
\end{equation}

The kinetic term is one half the mass times the squared length of the velocity
vector.  Differentiating and applying the multivariable chain rule gives

\begin{align}
  \frac{dE}{dt}
  &=m\sum_i\dot x_i\ddot x_i
    +\sum_i\frac{\partial V}{\partial x_i}\dot x_i
  \notag\\
  &=\sum_i\dot x_i
    \left(m\ddot x_i+\frac{\partial V}{\partial x_i}\right)
  =0.
  \label{eq:energy-derivative-multidimensional}
\end{align}

The proof is exactly the one-dimensional proof decorated with an index, but
its premise is stronger: the force components must all arise from one
potential.

\begin{figure}[htbp]
  \centering
  \includegraphics[width=0.84\textwidth]{lecture_02_figure_18.png}
  \caption{The indexed energy derivative.  Each bracket vanishes by the
  corresponding component of Newton's equation.}
  \lectureframe{01:28:30}
\end{figure}

For \(N\) particles in three dimensions, collect all coordinates into one list

\begin{equation}
  \begin{aligned}
    q_a={}&(x_1,y_1,z_1,x_2,y_2,z_2,\ldots,\\[-0.2em]
           &\phantom{(}x_N,y_N,z_N),\\
    &a=1,\ldots,3N.
  \end{aligned}
  \label{eq:flattened-coordinates}
\end{equation}

There are \(3N\) coordinate entries but only \(N\) particle masses; the same
mass occurs in the three entries belonging to a given particle.  If

\begin{equation}
  F_a=-\frac{\partial V(q_1,\ldots,q_{3N})}{\partial q_a},
  \label{eq:many-particle-conservative-force}
\end{equation}

then

\begin{equation}
  E=\sum_{A=1}^{N}\frac12m_A\lvert\dot{\mathbf r}_A\rvert^2
    +V(\mathbf r_1,\ldots,\mathbf r_N)
  \label{eq:many-particle-energy}
\end{equation}

is conserved by the same calculation.  Forces that arise from a potential are
called \emph{conservative} precisely because they support this conservation
law; the adjective has no political meaning.

\section{The Oscillator as an Energy Contour}
\enlargethispage{3\baselineskip}

For the spring force \(F=-kx\), Eq.~\eqref{eq:potential-definition-1d} gives

\begin{equation}
  V(x)=\frac12kx^2+V_0.
  \label{eq:spring-potential}
\end{equation}

Choosing \(V_0=0\), the oscillator energy is

\begin{equation}
  E=\frac12m\dot x^2+\frac12kx^2
   =\frac{p^2}{2m}+\frac12kx^2.
  \label{eq:oscillator-energy}
\end{equation}

For \(m=k=1\),

\begin{equation}
  2E=p^2+x^2=C^2.
  \label{eq:oscillator-energy-circle}
\end{equation}

The phase-space circles found earlier are therefore contours of constant
energy.  For a more general one-dimensional potential, the contours need not
be circles, but the motion still lies on curves

\begin{equation}
  \frac{p^2}{2m}+V(x)=E
  \label{eq:general-energy-contour}
\end{equation}

in the \((x,p)\) plane.  This returns us to the idea of the first lecture: when
trajectories remain on distinct cycles or orbits, a number can label those
orbits and remain unchanged throughout the motion.

We now have the state space, the Newtonian update equations, and the first two
conservation laws.  The next formulation of mechanics will begin from the
principle of least action and lead to Lagrangians, Hamiltonians, angular
momentum, and the more general structure behind these examples.

\begin{classroomqa}
  \audiencequestion{How closely do the distributed handout and the work on the
  board correspond, and where are the figures?\lecturetimestamp{01:37:56--01:38:54}}
  \lecturerresponse{The handout should follow the lecture closely; if material
  on the board is absent, the wrong file may have been sent.  The figures had
  not yet been prepared; they would appear, he jokes, when Continuing Studies
  paid the next check.}
\end{classroomqa}
